\section{Описание программы}

\subsection{Архитектура проекта}

Логика программы разделена на два модуля, что обеспечивает требование о чистых
функциях и позволяет тестировать их без ввода-вывода:

\begin{itemize}
  \item \texttt{array\_ops.cpp} реализует операции над динамическим массивом
        \texttt{int*} --- создание, изменение и освобождение, --- не обращаясь
        к \texttt{std::cin}; выводом занимается только \texttt{array\_print};
  \item \texttt{main.cpp} содержит текстовое меню, чтение и проверку ввода,
        обработку ошибок и вызов функций \texttt{array\_ops}.
\end{itemize}

Указатель и его длина хранятся в \texttt{main} как \texttt{int* arr} и
\texttt{std::size\_t size}. Операции, изменяющие длину, получают
\texttt{std::size\_t\& size} как out-параметр: функции
\texttt{array\_insert} и \texttt{array\_remove} обновляют длину по ссылке.
Функция поиска \texttt{array\_binary\_search} возвращает \texttt{bool}, а
найденный индекс передаётся через out-параметр
\texttt{std::size\_t\& out\_index}.

\subsection{Состав функций}

\begin{table}[htbp]
\centering
\caption{Чистые функции модуля \texttt{array\_ops}}
\begin{tabular}{@{}>{\raggedright\arraybackslash}p{0.96\textwidth}@{}}
\hline
\textbf{Сигнатура} \\
\hline
\texttt{int* array\_create(std::size\_t size)} \\
\texttt{void array\_delete(int*\& arr)} \\
\texttt{int* array\_resize(int* arr, std::size\_t size, std::size\_t new\_size)} \\
\texttt{int* array\_insert(int* arr, std::size\_t\& size, std::size\_t pos, int value)} \\
\texttt{int* array\_remove(int* arr, std::size\_t\& size, std::size\_t pos)} \\
\texttt{void array\_print(const int* arr, std::size\_t size)} \\
\texttt{void array\_selection\_sort(int* arr, std::size\_t size)} \\
\texttt{void array\_reverse(int* arr, std::size\_t size)} \\
\texttt{bool array\_binary\_search(const int* arr, std::size\_t size, int target,} \texttt{std::size\_t\& out\_index)} \\
\hline
\end{tabular}
\end{table}

\subsection{Алгоритм варианта}

Пункт~7 выполняет сортировку выбором по возрастанию, затем разворот массива
двумя указателями. Сортировка перебирает позиции \texttt{i} и находит индекс
минимума на отрезке \texttt{[i, size)}, меняя \texttt{arr[i]} и
\texttt{arr[minInd]} местами через временную переменную. Разворот сдвигает
указатели \texttt{left} и \texttt{right} к центру, обменивая
\texttt{arr[left]} и \texttt{arr[right]} до их встречи.

\subsection{Состав меню}

Меню соответствует семейству~A, дополнено пунктом~8 для демонстрации поиска:

\begin{table}[htbp]
\centering
\caption{Пункты меню}
\begin{tabular}{@{}r@{\quad}>{\raggedright\arraybackslash}p{4.3cm}@{\qquad}r@{\quad}>{\raggedright\arraybackslash}p{5.2cm}@{}}
\hline
1 & Создать массив & 5 & Изменить размер \\
2 & Напечатать & 6 & Печать \\
3 & Вставить элемент & 7 & Алгоритм варианта \\
4 & Удалить элемент & 8 & Бинарный поиск \\
  & & 0 & Выход \\
\hline
\end{tabular}
\end{table}

Пункт~8 является осознанным расширением меню семейства~A: цель работы прямо
требует «алгоритм варианта и поиск с возвратом результата и найденного значения
по ссылке», а обязательная функция семейства --- это
\texttt{array\_binary\_search}. Чтобы поиск выполнялся на упорядоченных данных,
пункт~8 создаёт копию массива, сортирует её выбором и выполняет бинарный поиск
по копии, не изменяя исходный массив. Копия освобождается сразу после поиска.

\subsection{Ввод, ошибки и владение памятью}

\begin{itemize}
  \item Размеры и позиции вводятся во временную переменную типа
        \texttt{long long} и проверяются на отрицательность до приведения к
        \texttt{std::size\_t}. Это исключает переполнение при преобразовании
        отрицательного числа в беззнаковый тип.
  \item Чтение выполняется через функции \texttt{readLongLong} и
        \texttt{readInt}, которые проверяют результат \texttt{std::cin} и при
        ошибке вызывают \texttt{clearInput()}: сбрасывают флаги состояния и
        \texttt{std::cin.ignore} до конца строки. Проверка
        \texttt{std::cin.eof()} при выборе пункта меню завершает программу
        штатно, поэтому перенаправление ввода из файла или пайпа не приводит к
        зацикливанию.
  \item Операции, выделяющие память (создание, вставка, удаление, изменение
        размера, копия для поиска), обёрнуты в \texttt{try/catch} по
        \texttt{std::bad\_alloc}. При нехватке памяти выводится сообщение, а
        исходный массив и его длина остаются неизменными.
  \item Каждая операция выводит результат: сам массив и новую длину. При
        вставке за конец выводится предупреждение и фактическая позиция
        вставки. При удалении позиции вне диапазона выводится допустимый
        диапазон, и массив не изменяется.
  \item Функции, изменяющие длину, работают по принципу «сначала новая память,
        затем копирование, затем \texttt{delete[]} старой». Поэтому при
        успешном выделении памяти исключение внутри функции уже не может
        привести к утечке. \texttt{array\_delete} обнуляет указатель.
\end{itemize}

\subsection{Сборка и запуск}

Проект собирается out-of-source через CMake: библиотека \texttt{array\_ops},
исполняемый файл \texttt{main} и тест \texttt{test\_variant}. Google Test
подключается через \texttt{FetchContent} (версия 1.14.0), тесты регистрируются
в CTest командой \texttt{gtest\_discover\_tests}.

\begin{lstlisting}[style=cmakestyle,caption={Сборка и запуск},numbers=none]
cmake -S . -B build
cmake --build build
./build/main
ctest --test-dir build --output-on-failure
\end{lstlisting}

\subsection{Исходный код}

\subsubsection{Система сборки}

\lstinputlisting[style=cmakestyle,caption={CMakeLists.txt}]{../../CMakeLists.txt}

\subsubsection{Заголовочный файл}

\lstinputlisting[style=cppstyle,caption={src/array\_ops.h}]{../../src/array_ops.h}

\subsubsection{Реализация операций}

\lstinputlisting[style=cppstyle,caption={src/array\_ops.cpp}]{../../src/array_ops.cpp}

\subsubsection{Точка входа}

\lstinputlisting[style=cppstyle,caption={src/main.cpp}]{../../src/main.cpp}

\subsubsection{Тесты}

\lstinputlisting[style=cppstyle,caption={tests/test\_variant.cpp}]{../../tests/test_variant.cpp}
